\documentclass[aspectratio=169]{beamer}
\usepackage[utf8]{inputenc}
\usepackage[T1]{fontenc}
\usepackage{amsmath,amssymb,booktabs}
\usepackage{hyperref}
\usetheme{Madrid}

\title{Norwegian Renewable Electricity Potential}
\subtitle{Hydropower, Wind, Solar, Grid Losses, Energy Efficiency, and System Balance}
\author{}
\date{}

\begin{document}

\frame{\titlepage}

\begin{frame}{Scope}
This short lecture summarizes six related questions for Norway:
\begin{enumerate}
    \item How much more hydropower can be obtained by upgrading existing facilities?
    \item What percentage increase in present electricity production could this represent?
    \item How much energy is currently lost in the grid, and what could be gained by reducing those losses?
    \item What are realistic uncertainty ranges for wind and solar additions?
    \item What are the main energy-efficiency and energy-saving potentials?
    \item How can one describe the full system balance between supply and demand?
\end{enumerate}
\end{frame}

\begin{frame}{Current Norwegian Electricity Production}
Recent Norwegian electricity production is approximately
\[
162\ \text{TWh/year}
\]
and mean annual hydropower production is about
\[
137\text{--}138\ \text{TWh/year}.
\]

\medskip

Current wind production is about
\[
15.4\ \text{TWh/year}.
\]
\end{frame}

\begin{frame}{Potential from Upgrading Existing Hydropower}
A widely cited estimate is that upgrading and expanding existing hydropower plants could provide at least
\[
7\ \text{TWh/year}
\]
without major new environmental conflicts.

\medskip

This corresponds to:
\[
\frac{7}{137}\approx 5.1\%
\]
of present hydropower production, and
\[
\frac{7}{162}\approx 4.3\%
\]
of current total electricity production.
\end{frame}

\begin{frame}{Grid Losses in the Norwegian Power System}
A recent Statnett figure gives consumption of
\[
139.2\ \text{TWh}
\]
including grid losses, versus
\[
131.5\ \text{TWh}
\]
without grid losses.

\medskip

Hence the implied losses are
\[
139.2-131.5 = 7.7\ \text{TWh/year}.
\]

This corresponds approximately to
\[
5.5\% \text{ of gross consumption}
\]
and
\[
4.8\% \text{ of total production.}
\]
\end{frame}

\begin{frame}{What Could Be Gained by Reducing Grid Losses?}
Illustrative scenarios from the current loss level \(7.7\ \text{TWh/year}\):
\[
10\% \text{ reduction} \Rightarrow 0.77\ \text{TWh/year}
\]
\[
25\% \text{ reduction} \Rightarrow 1.93\ \text{TWh/year}
\]
\[
50\% \text{ reduction} \Rightarrow 3.85\ \text{TWh/year}
\]

\medskip

As shares of current production (\(\sim 162\ \text{TWh}\)):
\[
0.5\%,\quad 1.2\%,\quad 2.4\%.
\]
\end{frame}

\begin{frame}{Wind Power: Current Level and Long-Term Policy}
Current wind production is approximately
\[
15.4\ \text{TWh/year}.
\]

\medskip

The government's long-term offshore ambition is
\[
30\ \text{GW by 2040},
\]
which corresponds to a conservative estimate of
\[
120\ \text{TWh/year}.
\]

\medskip

This is a very large technical-policy ambition, not a short-term forecast.
\end{frame}

\begin{frame}{Wind Power: Uncertainty Ranges}
A useful way to think about wind is to distinguish three levels:

\begin{itemize}
    \item \textbf{Current realized level:} \(\sim 15\ \text{TWh/year}\)
    \item \textbf{Medium-term incremental buildout:} uncertain, depends on licensing, grid, and industrial demand
    \item \textbf{Long-term offshore policy ambition:} up to \(\sim 120\ \text{TWh/year}\)
\end{itemize}

\medskip

Thus the uncertainty range is very wide:
\[
\text{small additions in the 2020s} \quad \text{vs.} \quad \text{very large additions by 2040.}
\]
\end{frame}

\begin{frame}{Solar Power: Policy Target and Uncertainty}
Norway has a political target for new solar power of
\[
8\ \text{TWh by 2030}.
\]

\medskip

Solar remains highly uncertain because output depends on:
\begin{itemize}
    \item rooftop and distributed deployment,
    \item local grid access,
    \item investment incentives,
    \item seasonal mismatch between summer output and winter demand.
\end{itemize}

\medskip

A reasonable uncertainty framing is:
\[
\text{low realization} < 8\ \text{TWh}, \qquad
\text{policy success} \approx 8\ \text{TWh by 2030}.
\]
\end{frame}

\begin{frame}{Energy Efficiency: Official Targets}
Norway has a target to reduce energy use in existing buildings by
\[
10\ \text{TWh by 2030}
\]
relative to 2015 levels.

\medskip

Norway also has a national target to improve energy intensity in the mainland economy by
\[
30\%
\]
from 2015 to 2030.
\end{frame}

\begin{frame}{Energy Saving Technologies}
Examples of technologies and measures that reduce consumption:
\begin{itemize}
    \item better insulation and building envelopes,
    \item heat pumps and smarter control systems,
    \item better ventilation and heat recovery,
    \item more efficient industrial motors and process control,
    \item smart charging, flexible demand, and load shifting,
    \item more efficient household appliances.
\end{itemize}

\medskip

NVE has also estimated that ecodesign and energy labelling for household appliances can save around
\[
3\ \text{TWh in 2030}.
\]
\end{frame}

\begin{frame}{Uncertainty Ranges for Energy Efficiency}
A useful decomposition is:
\begin{itemize}
    \item \textbf{Buildings:} around \(10\ \text{TWh}\) policy target by 2030
    \item \textbf{Household appliances:} around \(3\ \text{TWh}\) by 2030
    \item \textbf{Additional industry and flexibility gains:} uncertain and highly scenario-dependent
\end{itemize}

\medskip

Hence a broad but defensible range for total energy-saving effects by 2030 is:
\[
10\text{--}20\ \text{TWh/year},
\]
with a central anchor near the lower half of that interval.
\end{frame}

\begin{frame}{Supply-Side Uncertainty Summary}
\begin{center}
\begin{tabular}{lcc}
\toprule
Measure & Central figure & Interpretation \\
\midrule
Hydropower upgrades & \(\sim 7\) TWh & relatively firm \\
Current wind & \(\sim 15.4\) TWh & realized production \\
Offshore wind 2040 & \(\sim 120\) TWh & long-term ambition \\
Solar 2030 target & \(\sim 8\) TWh & policy target \\
Grid-loss recovery & \(< 7.7\) TWh & technical upper bound \\
\bottomrule
\end{tabular}
\end{center}
\end{frame}

\begin{frame}{Demand-Side Uncertainty Summary}
Demand is increasingly shaped by:
\begin{itemize}
    \item electrification of transport,
    \item data centres,
    \item industrial decarbonization,
    \item hydrogen and battery projects,
    \item energy-efficiency improvements.
\end{itemize}

\medskip

Statnett's long-term market analysis gives a wide range for Norwegian electricity consumption in 2050:
\[
180\text{--}260\ \text{TWh}.
\]

The low scenario assumes more efficiency and little new production; the high scenario assumes large new production and stronger industrial growth.
\end{frame}

\begin{frame}{A Full System Model: Supply Minus Demand}
A simple system-balance identity is
\[
B = P_{\mathrm{hydro}} + P_{\mathrm{wind}} + P_{\mathrm{solar}} + P_{\mathrm{other}}
+ \Delta P_{\mathrm{upgrade}}
+ \Delta P_{\mathrm{loss\ reduction}}
- D_{\mathrm{gross}}
+ \Delta D_{\mathrm{efficiency}},
\]
where:
\begin{itemize}
    \item \(B\) is the energy balance,
    \item production terms enter positively,
    \item demand enters negatively,
    \item efficiency reduces demand and therefore improves balance.
\end{itemize}
\end{frame}

\begin{frame}{Baseline 2025 System Balance}
A simple baseline picture is:
\[
P_{\mathrm{total}} \approx 162\ \text{TWh},
\qquad
D_{\mathrm{gross}} \approx 139.2\ \text{TWh}.
\]

Hence the current gross surplus is roughly
\[
162 - 139.2 \approx 22.8\ \text{TWh}.
\]

\medskip

This does \emph{not} mean that the system is unconstrained:
\begin{itemize}
    \item regional bottlenecks matter,
    \item dry years matter,
    \item future load growth can quickly absorb the surplus.
\end{itemize}
\end{frame}

\begin{frame}{Illustrative 2030 Balance Scenarios}
One can write a stylized balance as:
\[
B_{2030} \approx 162 + \Delta P_{\mathrm{hydro}} + \Delta P_{\mathrm{wind}} + \Delta P_{\mathrm{solar}} + \Delta P_{\mathrm{loss}} - \Delta D.
\]

Illustrative cases:
\begin{itemize}
    \item \textbf{Efficiency-led case:} moderate new production, stronger demand restraint
    \item \textbf{Production-led case:} more wind and solar, faster electrification, weaker demand restraint
    \item \textbf{Tight-balance case:} demand grows faster than new production
\end{itemize}

The system outcome is therefore highly sensitive to both industrial demand growth and the pace of new generation.
\end{frame}

\begin{frame}{Policy Scenarios and Tension in the Outlook}
Two broad messages coexist:
\begin{itemize}
    \item Statnett's long-term scenarios assume balanced growth over time, with demand in 2050 between \(180\) and \(260\) TWh.
    \item DNV warns of a short-term electricity supply-demand deficit in the early 2030s unless new production and efficiency measures arrive fast enough.
\end{itemize}

\medskip

So the main risk is not the ultimate technical resource base, but the timing mismatch between:
\[
\text{new demand} \quad \text{and} \quad \text{new supply + savings}.
\]
\end{frame}

\begin{frame}{Comparison of Levers}
\begin{center}
\begin{tabular}{lcc}
\toprule
Lever & Typical scale & Horizon \\
\midrule
Hydropower upgrades & \(\sim 7\) TWh & near / medium term \\
Grid-loss reduction & \(1\text{--}4\) TWh & gradual \\
Solar target & up to \(8\) TWh & by 2030 \\
Efficiency & \(10\text{--}20\) TWh & by 2030, uncertain \\
Offshore wind ambition & up to \(120\) TWh & by 2040 \\
\bottomrule
\end{tabular}
\end{center}
\end{frame}

\begin{frame}{Main Conclusion}
A reasonable integrated conclusion is:
\begin{itemize}
    \item Existing hydropower upgrades are important but limited in scale (\(\sim 7\) TWh).
    \item Grid-loss reduction helps, but is smaller than the main production and efficiency levers.
    \item Solar is still small today, but policy aims for up to \(8\) TWh by 2030.
    \item Wind has the largest long-term technical upside, especially offshore.
    \item Energy efficiency is a major \emph{virtual resource}, plausibly in the \(10\text{--}20\) TWh range by 2030.
    \item The key Norwegian challenge is managing the timing of demand growth relative to new supply and efficiency gains.
\end{itemize}
\end{frame}

\begin{frame}{References}
\footnotesize
\begin{itemize}
    \item Statnett, Annual Report 2025.
    \item Statnett, Long-term Market Analysis 2024--2050.
    \item SINTEF, ``O/U -- Upgrading and expanding hydropower''.
    \item The Norwegian Offshore Directorate, energy production fact box.
    \item IEA, Norway 2022 executive summary.
    \item NVE, ecodesign and energy labelling savings estimate.
    \item Government of Norway, Roadmap 2.0 The Green Industrial Initiative.
    \item DNV, Energy Transition Outlook Norway 2025.
\end{itemize}
\end{frame}

\end{document}
